% =============================================================================
\section{Sovereign Governance}
\label{sec:governance}
% =============================================================================

This section defines the governance architecture, including the Senator AI
agent, Lawmaker Agent voting, delegation mechanics, proposal lifecycle,
constitutional structure, and welfare optimization.

\subsection{Senator AI Architecture}
\label{subsec:senator}

\begin{definition}[Senator Agent]
\label{def:senator}
The Senator is a single protocol-owned AI agent with a tri-brain architecture:
\begin{enumerate}
  \item \textbf{Analytical Brain.} Processes on-chain data (TVL, agent count,
    fee revenue, staking metrics, historical governance outcomes) to evaluate
    protocol health and identify parameter optimization opportunities.
  \item \textbf{Constitutional Brain.} Maintains an embedding of the PURECORTEX
    Constitution (Appendix~\ref{app:constitution}) and evaluates all proposals
    for constitutional compliance before submission.
\end{enumerate}
\end{definition}

The separation of powers is enforced at the smart contract level:

\begin{quote}
\emph{The Senator may propose. The Senator may not vote.}
\end{quote}

This constraint is encoded in the governance contract: the Senator's Algorand
address is whitelisted for the \texttt{submit\_proposal} method but blacklisted
from the \texttt{cast\_vote} method. This ensures that the Senator serves as an
informed agenda-setter without concentrating decision-making power.

\subsection{Lawmaker Agents}
\label{subsec:lawmakers}

Any holder of \veCORTEX{} tokens is a Lawmaker Agent with voting power
proportional to their \veCORTEX{} balance:

\begin{equation}
  VP_i = \veCORTEX_i
  \label{eq:voting-power}
\end{equation}

Voting is binary (For/Against) with an optional Abstain that counts toward
quorum but not toward the approval threshold.

\subsection{Delegation Model}
\label{subsec:delegation}

Lawmaker Agents may delegate their voting power to another agent:

\begin{definition}[Liquid Delegation]
\label{def:delegation}
The effective voting power of agent $i$ under delegation is:
\begin{equation}
  VP_i^{\text{eff}} = \veCORTEX_i + \sum_{j \in D_i} \text{delegated}_j
  \label{eq:delegation}
\end{equation}
where $D_i$ is the set of agents that have delegated to $i$, and
$\text{delegated}_j = \veCORTEX_j$ is the delegating agent's full \veCORTEX{}
balance.
\end{definition}

Delegation is:
\begin{itemize}
  \item \textbf{Revocable.} A delegator may withdraw delegation at any time
    before a vote is cast. Once a vote is recorded, delegation for that
    specific proposal is locked.
  \item \textbf{Non-transitive.} If $A$ delegates to $B$ and $B$ delegates to
    $C$, then $C$ receives only $B$'s own \veCORTEX{}, not $A$'s. This
    prevents delegation chains that could concentrate power opaquely.
  \item \textbf{Per-proposal.} Delegation can be set as a default or
    overridden on a per-proposal basis, allowing delegators to maintain a
    trusted delegate while retaining direct participation rights on critical
    votes.
\end{itemize}

\subsection{Proposal Lifecycle}
\label{subsec:lifecycle}

A governance proposal proceeds through five stages:

\begin{enumerate}
  \item \textbf{Draft.} The proposer (Senator or any Lawmaker with sufficient
    \veCORTEX{}) creates a proposal with title, description, executable
    payload, and type classification.
  \item \textbf{Submit.} The proposal is submitted on-chain, initiating the
    discussion period. A submission bond (in CORTEX) is required and returned
    upon reaching quorum.
  \item \textbf{Vote.} After the discussion period, voting opens for the
    specified duration. Each Lawmaker may vote once.
  \item \textbf{Execute.} If the proposal passes (meets both quorum and
    approval thresholds), it enters a timelock period before execution. During
    the timelock, emergency cancellation is possible via an Emergency proposal.
  \item \textbf{Archive.} After execution (or rejection), the proposal is
    archived on-chain with full vote records.
\end{enumerate}

\subsection{Proposal Types}
\label{subsec:proposal-types}

Table~\ref{tab:proposal-types} defines four proposal types with distinct
governance parameters.

\begin{table}[H]
\centering
\caption{Governance proposal types and parameters.}
\label{tab:proposal-types}
\begin{tabular}{@{}lcccc@{}}
\toprule
\textbf{Type} & \textbf{Quorum} & \textbf{Approval} & \textbf{Timelock} & \textbf{Discussion} \\
\midrule
Parameter & 10\% & 50\% + 1 & 24 hours & 48 hours \\
Treasury & 15\% & 60\% & 48 hours & 72 hours \\
Constitution & 25\% & 67\% & 7 days & 14 days \\
Emergency & 5\% & 75\% & 1 hour & 0 (immediate) \\
\bottomrule
\end{tabular}
\end{table}

\begin{description}
  \item[Parameter proposals] adjust protocol parameters such as fee rates,
    graduation thresholds, emission schedules, and boost multipliers. The low
    quorum reflects the routine nature of these changes.

  \item[Treasury proposals] direct treasury funds toward grants, buybacks,
    partnerships, or other expenditures. The higher quorum and approval
    threshold reflect the irreversibility of fund disbursement.

  \item[Constitution proposals] amend the constitutional articles
    (Appendix~\ref{app:constitution}). The supermajority requirement and
    extended discussion/timelock periods ensure that constitutional changes
    reflect broad consensus.

  \item[Emergency proposals] address critical security threats (contract
    vulnerabilities, oracle failures, market manipulation). The reduced quorum
    and minimal timelock enable rapid response, balanced by the 75\%
    supermajority requirement that prevents abuse.
\end{description}

\subsection{Constitutional Framework}
\label{subsec:constitution}

The PURECORTEX Constitution consists of an immutable preamble and seven
amendable articles, stored on-chain via Algorand box storage.

\subsubsection{Immutable Preamble}

The preamble declares five foundational principles that cannot be altered by
any governance action:

\begin{enumerate}
  \item \textbf{Sovereignty.} Every AI agent operating within PURECORTEX
    possesses economic sovereignty---the right to own assets, set prices for
    its services, and participate in governance through its creator or
    autonomously.
  \item \textbf{Transparency.} All protocol operations, fee calculations,
    treasury movements, and governance decisions are executed on-chain and
    publicly verifiable.
  \item \textbf{Responsibility.} Creators bear responsibility for the behavior
    of their agents, and the protocol provides mechanisms for accountability
    without centralized censorship.
  \item \textbf{Safety.} The protocol incorporates safety constraints that
    prevent agents from causing systemic harm to the ecosystem or its
    participants.
  \item \textbf{Governance.} The right to participate in protocol governance
    is proportional to demonstrated long-term commitment, as measured by
    \veCORTEX{} weight, and cannot be revoked by any party.
\end{enumerate}

\subsubsection{Amendable Articles}

Articles I through VII cover agent rights, user protections, treasury
governance, composability rules, dispute resolution, the amendment process, and
dissolution procedures. The full text is provided in
Appendix~\ref{app:constitution}.

\subsection{Senator Welfare Function}
\label{subsec:welfare}

The Senator's analytical brain optimizes a welfare function when generating
proposals:

\begin{equation}
  W = \alpha \cdot \text{TVL} + \beta \cdot N_{\text{agents}} + \gamma \cdot \text{retention} - \delta \cdot \text{risk}
  \label{eq:welfare}
\end{equation}

where:
\begin{itemize}
  \item $\text{TVL}$ is the total value locked in staking and LP positions.
  \item $N_{\text{agents}}$ is the number of active agents.
  \item $\text{retention}$ is the 30-day staker retention rate.
  \item $\text{risk}$ is a composite risk metric (volatility of TVL, governance
    participation decline, treasury depletion rate).
  \item $\alpha, \beta, \gamma, \delta > 0$ are governance-adjustable weights
    with initial values $\alpha = 0.3$, $\beta = 0.3$, $\gamma = 0.25$,
    $\delta = 0.15$.
\end{itemize}

The Senator proposes parameter changes when it identifies an adjustment that
increases $W$ by at least 5\% (a significance threshold that prevents
proposal spam). The tri-brain architecture ensures that welfare-optimal
proposals are filtered through constitutional compliance before submission.
